\noindent Newcomers may build connections with each other, or with already established individuals. Ties may take many forms. For example, we might rely on someone for advice but not for lending us money. A conceptual framework to study homophily in newcomers' networks must take this into account, and marry several aspects of tie formation validated by the literature. Note that throughout, I will refer to \textit{the institution} as whatever gathers the newcomers in one place and settles who interacts with whom. Here, it is the industrial park that recruits, trains and assigns them to production units.  \\

\parhead{Homophily drivers}

\noindent Firstly, before asking why a tie between two people took the form it did, one must ask whether that tie was made after the two arrived. A tie either predates the institution or was made inside it. People might emigrate with many close relationships, and essentially stick with them throughout because they do not need any other. I will dub this the \textit{importation channel}. If this channel dominates, the composition of the newcomer's network is settled before the institution mixes anyone. This does not say much about how ties form after arrival, it only isolates the portion of homophily attributable to importation. The rest is to be explained by other mechanisms. \\

\noindent Conditional on a tie having been formed at the institution, three mechanisms may be in action. The first one is the \textit{opportunity channel}. If newcomers meet people only through the institution, and the institution assembles newcomers by origin, then the newcomers' network will be homophilous. The second, and most obvious one, is the \textit{taste channel}. If newcomers meet everyone but prefer to be with their own, their network will  be homophilous. The third is the \textit{exclusion channel}. Assuming newcomers seek ties outside of their own group, but that they are refused, then their network will also end up to be homophilous. \\

\parhead{Empirical evidence}

\noindent Each of the four channels has been documented and validated, though not in the same setting. \\

\noindent That newcomers' networks are \textit{imported} is established, although the clearest evidence comes from papers on how firms hire.  If a firm recruits through its incumbents, a share of every intake arrives already tied to someone inside. This is often the case. Bangladeshi garment factories hire through referrals deliberately, because a referral lets the firm hold the referring worker liable for the recruit's output \citep{heath2018why}. Around 30\% percent of workers in her sample received a referral in their current job, 65\% of those referrals came from a relative, and 60\% of referred workers began on the same production team as the person who referred them. In Danish administrative data, a tie raises the probability of being hired at a firm by 74\%, of which only about a tenth comes from a higher success rate per application and all the rest from where workers choose to apply \citep{harmon2025coworker}. The same holds from the worker's side, and for migrants specifically. Among millions of Chinese domestic migrants on a food delivery platform, 60\%  of new workers join with a referrer already working on the platform, and two thirds of those pairs share a county of birth \citep{tang2025follow}. Note what this implies for measurement - the hiring cohort is itself a network outcome. \\


\noindent That \textit{opportunity} drives tie formation is even better documented. It is measured at scale by \citet{chetty2022social}. They document a gap in cross-class friendship by analysing the friendship links of 70.3 million Facebook users aged 25 to 44 (up to 82\% of that age group). They find that exposure drives about half of that shortfall. Meetings between people of different class have essentially little chance of happening, given that low-status people tend to belong to institutions containing no high-status people\footnote{In the paper, they refer to schools, workplaces and religious groups.}. The other half of the mentioned shortfall is friending bias, \textit{i.e.} the tendency of low-status people to befriend high-status people at lower rates conditional on exposure. It appears that cross-group contact generates ties only if it is collaborative and prolonged. An experiment run in India found that being placed in the same cricket team raises cross-caste friendship and lowers own-caste favouritism, while being placed in adversarial teams does neither \citep{lowe2021types}. \\

\noindent Evidence on \textit{taste} is somewhat thinner, mostly because taste is usually thought of as what remains once opportunity has been controlled for. The exceptions are designs in which the choice set is fixed and every decision is recorded. Same-race preferences emerge in a speed dating experiment on Columbia graduate students, and the strength of these preferences is predicted by the racial composition of the ZIP code in which the subject grew up \citep{fisman2008racial}. In friendships and study partnerships among university students, homophily on gender and ethnicity appears early and stays constant over time \citep{jackson2022dynamics}. Note that separating opportunity from taste is an old problem, with \citet{currarini2009economic} showing that the same segregation patterns can be generated either by biases in preferences or by biases in meetings. \citet{chetty2022social} encounters the same difficulty. They argue that the previously discussed friending bias is not a preference parameter, as it may reflect biased sampling from the pool of available peers. \\

\noindent \citet{chetty2022social} also stress that their friending bias measure cannot capture what I call \textit{exclusion}. Given that tie requires the consent of both parties \citep{jackson1996strategic}, a newcomer with few out-group ties may have been refused rather than uninterested (\textit{i.e.,} we cannot distinguish somebody who refused a tie from someone who was never approached). An unreciprocicated tie might be evidence of an unbalance of power in relationships. Using a census of 119 households in a Tanzanian village, \citet{comola2014testing} test what a self-reported tie actually captures, and find that responses are best explained by a willingness to link rather than a realised relationship.  They also find evidence of self-censoring, when the respondent omits an individual from her list not because she does not have a tie with that person, but because she anticipates that the other person would refuse her. More generally, exclusion has economic consequences. It has been found that physicians refer disproportionately to specialists of their own gender, a bias driven by the referring physician's decision rather than by sorting of patients or of physicians, and because most referrers are male the net effect is a 5\% lower demand for female specialists \citep{zeltzer2020gender}. \\

\noindent One caveat cuts across all four channels. A channel may act on which functions a tie carries, rather than on whether it exists at all. \citet{chandrasekhar2025multiplexing} map socialising, advising, helping and lending as separate layers across Indian villages, and find that most layer pairs correlate above 0.5. This implies that the same people recur across functions, but no one layer stands in for another. The variables ordinarily substituted for network data (\textit{e.g.} ethnicity or geographic proximity) are not good proxies and overpredict ties \citep{conley2010learning}. 
